% TOP_PROBLEM: {"id": 30002148, "title": "Regular Conformal Equations for Matter with Nonzero Trace", "category_id": 16, "category": {"id": 16, "name": "physics", "display_name": "Mathematical Physics", "description": "Problems at the intersection of mathematics and physics.", "slug": "physics", "order_index": 16, "created_at": "2026-07-31T15:26:25.671Z"}, "status": "open", "statusQualification": "The question is a construction programme across matter models. It does not say that one fixed rescaling must work for every imaginable nonzero-trace tensor; admissible matter equations must be specified as in the source."}
% FIRST_POSTED: 2026-10-11
% ENTRY_KIND: statement_only
% SUBMISSION: {"kind": "statement_only", "category": "top_problems", "problem_id": "30002148", "proposal_type": "established", "source": "Mathematical Aspects of General Relativity (2012), Oberwolfach Reports, DOI 10.4171/owr/2012/37, pp. 2269--2333. Primary report accessed 2026-10-11; relevant PDF pages 47. https://doi.org/10.4171/owr/2012/37; Friedrich et al., conformal methods review (catalogue-listed context; not independently consulted in this preparation). https://pmc.ncbi.nlm.nih.gov/articles/PMC10788156/", "status_evidence": "Canonical source status open retained. Supplied literature triage dated 2026-08-21: Later conformal work handles trace-free matter and selected non-trace-free sources, but no general regular conformal system for arbitrary nonzero-trace stress tensors has been established. Primary-report access checked 2026-10-11; this historical report and inherited triage do not certify current resolution of every component.", "status_checked_on": "2026-10-11", "duplicate_check": "Compared canonical ID, title and normalized target with the complete local UnsolvedMath registry; normalized exact-text matches: none; title matches: none. Related or overlapping canonical records: none identified in this batch. Primary source context and the existing OWR statement audit used for scope and duplicate review; an exhaustive mathematical-equivalence/literature search was not performed.", "attribution": "Original proposer not named in the supplied canonical record; attribution remains to the identified Oberwolfach contribution/report, not its editors as invented proposers. Explanatory statement preparation and source/semantic audit assisted by OpenAI GPT-6 Codex; no human mathematical authorship or independent proof verification is claimed."}
% SOURCE_STATEMENT: "For matter with a nonzero-trace energy-momentum tensor, does there exist a choice of conformal variables yielding a regular conformal field-equation system with matter?"
% Preparation and final self-audit: CONTRIBUTING.md prompts C and D; source/audit artifacts retained in the local physics_statements_2026_10_11 batch.
% Imported from: physics_statements.tex; UnsolvedMath ID 30002148
% !TeX program = lualatex
\documentclass[11pt]{article}
\usepackage[margin=25mm]{geometry}
\usepackage{fontspec}
\setmainfont{Latin Modern Roman}
\usepackage{amsmath,amssymb,mathtools}
\usepackage{unicode-math}
\setmathfont{Latin Modern Math}
\usepackage{xurl,hyperref}
\hypersetup{colorlinks=true,urlcolor=blue}
\setcounter{secnumdepth}{0}
\setlength{\parindent}{0pt}
\setlength{\parskip}{5pt}
\setlength{\emergencystretch}{3em}
\newcommand{\sref}[2]{\hyperlink{source:#1:#2}{[S#2]}}
\newcommand{\cataloguescope}[1]{\paragraph{Recorded target.}#1}
\newcommand{\R}{\mathbb R}
\newcommand{\C}{\mathbb C}
\newcommand{\N}{\mathbb N}
\newcommand{\Z}{\mathbb Z}
\newcommand{\Q}{\mathbb Q}
\newcommand{\F}{\mathbb F}
\newcommand{\eps}{\varepsilon}
\begin{document}
% UnsolvedMath ID 30002148; source catalogue code OWR-12009-005
\setcounter{section}{30002147}
\section{Regular Conformal Equations for Matter with Nonzero Trace}\label{30002148}
{\small UnsolvedMath ID 30002148 · Mathematical Physics}
\subsection{Definitions and mathematical statement}

\paragraph{Shared notation.}$\mathbb N=\{1,2,\ldots\}$; the standard number fields have their usual meanings. All additional conventions are stated locally.


\paragraph{Definitions and conventions.}
A conformal compactification writes the physical metric as $\widetilde g=\Omega^{-2}g$, with a conformal boundary at $\Omega=0$. The trace of a physical energy--momentum tensor is $\widetilde g^{ab}T_{ab}$. A regular conformal system is an equivalent system in rescaled metric, curvature and matter variables whose coefficients and unknowns do not acquire singular divisions by $\Omega$ at the boundary. Matter equations and their transformation law are essential parts of such a system.

\paragraph{Target retained from the canonical source catalogue.}
\paragraph{Catalogue wording retained for traceability.}
For matter with a nonzero-trace energy-momentum tensor, does there exist a choice of conformal variables yielding a regular conformal field-equation system with matter?

\paragraph{Source scope and qualification.}
The question is a construction programme across matter models. It does not say that one fixed rescaling must work for every imaginable nonzero-trace tensor; admissible matter equations must be specified as in the source.
\subsection{Short English statement}
Can conformal field equations remain regular at the conformal boundary when the matter stress tensor has nonzero trace?
\subsection{Sources}
\begin{itemize}
\item[\hypertarget{source:30002148:1}{S1}] ULAM.AI and the UnsolvedMath Contributors, UnsolvedMath, record 30002148 (OWR-12009-005). Catalogue curation under CC BY 4.0; definitions and English exposition added locally.
\url{https://huggingface.co/datasets/ulamai/UnsolvedMath}.
\item[\hypertarget{source:30002148:2}{S2}] Mathematical Aspects of General Relativity (2012), Oberwolfach Reports, DOI 10.4171/owr/2012/37, pp. 2269--2333. Primary report accessed 2026-10-11; relevant PDF pages 47.
\url{https://doi.org/10.4171/owr/2012/37}.
\item[\hypertarget{source:30002148:3}{S3}] Friedrich et al., conformal methods review (catalogue-listed context; not independently consulted in this preparation).
\url{https://pmc.ncbi.nlm.nih.gov/articles/PMC10788156/}.
\end{itemize}
\label{end:30002148}
\end{document}
